\documentclass[8pt]{strategynet-journal}
\usepackage{volt-paper-style}
\hypersetup{pdftitle={Pricing and Risk Analytics with Volt: Applications and APIs},pdfauthor={research@strategynet.ai}}
\newcommand{\VOLT}{\textsf{Volt}}
\newcommand{\E}{\mathbb{E}}
\SNHeader{strategynet.ai:Volt}
\title{Pricing and Risk Analytics with Volt\\\large Applications, APIs and CPU/GPU Execution}
\author{\href{mailto:research@strategynet.ai}{research@strategynet.ai}}
\date{October 2026 --- Tutorial and implementation report}
\begin{document}
\maketitle
\begin{abstract}
\normalsize An options book can contain thousands of contracts whose values and hedge requirements change with each market update. A risk application needs to connect those observations to contract prices, portfolio exposures and the effects of possible stock or volatility moves. Volt provides batch valuation and sensitivity APIs, a quote-to-surface risk service and a persistent learner interface for hedging applications. We follow a small call-and-put book through those tools, explain the units returned by each calculation, and show how automatic differentiation and CPU/GPU execution support larger workloads. On an RTX 4090 test system, valuing 65,536 rows already held on the device takes 0.67 milliseconds on CUDA and 1.75 on CPU; a complete first-order Greek report with spot gamma takes 27.66 milliseconds on CUDA and 24.24 on CPU. The resident valuation benefits from CUDA; the complete Greek request runs faster on CPU at this size. Independent numerical controls cover Black--Scholes, shared-factor sensitivities, Monte Carlo estimators and Heston pricing.
\end{abstract}
\begin{multicols}{2}
\tableofcontents
\end{multicols}
\thispagestyle{fancy}
\clearpage

\begin{SNcolumns}
\section{Pricing and risk in an application}
Consider a desk holding calls and puts on a stock. A new market observation
changes the estimated value of those contracts and the positions needed to
hedge them. The desk wants both numbers quickly. It also needs to understand
what caused the change: a move in the stock price, a change in volatility, the
passage of time, or a change in the model's assumptions.

For a call with delta 0.55, a one-unit stock rise increases its value by about
0.55 per underlying unit. A trader who has sold the call can buy 0.55 shares
to offset that first-order exposure. As the stock moves, delta changes; gamma
measures that response. Vega describes the change in option value as the
assumed volatility moves. These sensitivities form the book's Greeks.

For 256 risk inputs, central differences require 512 portfolio repricings:
one upward and one downward move for each input. Every added contract then
adds work to each of those repricings. Automatic differentiation shares the
price calculation across its input sensitivities, while a GPU can evaluate
many contract rows in parallel. Volt exposes these calculations through
interfaces that a batch process, a risk panel or a hedging application can call.


\end{SNcolumns}
\begin{center}
\begin{minipage}{\textwidth}
\noindent Figure~\ref{fig:execution} shows how the application request reaches its
financial and numerical calculation, and where the resulting tensors live.
\par\medskip\centering
\begin{tikzpicture}[>=Latex,font=\normalsize,
 box/.style={draw=darkblue,rounded corners=2pt,align=center,inner sep=5pt,minimum height=10mm},
 compute/.style={box,fill=lightblue!6},
 memory/.style={box,fill=darkgreen!5}]
 \node[box,text width=2.2cm] (input) at(0,0) {Quotes and\\positions};
 \node[compute,text width=2.8cm] (volt) at(3.2,0) {Volt APIs\\Contract inputs\\Units and models};
 \node[compute,text width=2.8cm] (burn) at(7.0,0) {Burn\\Tensors and\\differentiation};
 \node[memory,text width=2.8cm] (cpu) at(11.2,0) {CPU backend\\Host RAM};
 \node[compute,text width=2.8cm] (scalar) at(3.2,-2.0) {Live quote API\\CPU IV and surface\\Contract Greeks};
 \node[compute,text width=2.8cm] (cube) at(7.0,-2.0) {CubeCL\\GPU kernels and runtime};
 \node[memory,text width=2.8cm] (metal) at(11.2,-1.75) {Metal GPU\\Shared system RAM};
 \node[memory,text width=2.8cm] (cuda) at(11.2,-3.25) {CUDA GPU\\Device memory\\(VRAM)};
 \node[box,text width=2.2cm] (hyper) at(0,-4.6) {Hypercube\\Published risk snapshots};
 \node[compute,text width=2.8cm] (results) at(3.2,-4.6) {API results\\Contract marks\\Portfolio risk};
 \draw[->] (input)--(volt);\draw[->] (volt)--(burn);\draw[->] (burn)--(cpu);
 \draw[->] (volt)--(scalar);\draw[->] (scalar)--(results);
 \draw[->] (burn)--(cube);\draw[->] (cube)--(metal);\draw[->] (cube)--(cuda);
 \draw[darkgreen] (cpu.east)--(13.3,0)--(13.3,-4.6);
 \draw[darkgreen] (metal.east)--(13.3,-1.75);
 \draw[darkgreen] (cuda.east)--(13.3,-3.25);
 \draw[->,darkgreen] (13.3,-4.6)--node[above]{Results returned to host}(results.east);
 \draw[->] (results)--(hyper);
\end{tikzpicture}
\captionof{figure}{Volt connects the financial request to its numerical route.
Burn manages tensors and derivatives; CubeCL supplies GPU execution. CPU data
uses host RAM, Apple Metal resources use shared system RAM, and CUDA tensors
use GPU memory. Completed portfolio results are published through Hypercube.}
\label{fig:execution}
\end{minipage}
\end{center}

\begin{SNcolumns}
\section{The request and result APIs}
The application submits a contract book and an as-of market snapshot. The book
supplies contract identity, signed quantity and multiplier; the snapshot
supplies the underlying price, volatility, rates, dividend yield and observation
time. The \texttt{VoltValuationRequest} interface checks those records and
projects supported contracts into the numerical calculation~\cite{volt}.
The current valuation adapter accepts unexpired European calls and puts.
Their exercise and settlement records remain attached to the contract.

Applications can use the native library calls or the supplied command-line
runners. A valuation result contains contract marks and an aggregated book
value. The sensitivity API adds delta, vega, rho, calendar theta and dividend-yield
sensitivity; callers can request spot gamma as well. Per-contract outputs
retain the input row identity, and the portfolio calculation applies the
original signed quantities and multipliers. This lets a risk panel trace a
large exposure back to its contributing positions.

\begin{center}
\begin{tabular}{@{}>{\raggedright\arraybackslash}p{2.0cm}p{5.2cm}@{}}
\toprule
Application & Supplied calculation\\
\midrule
Batch risk & Contract values, Greeks and signed portfolio totals\\
Scenario analysis & Revaluation under supported stock and volatility changes\\
Live option view & Quote-implied volatilities, fitted surface, book Greeks and scenarios\\
Hedge learning & Continuous holdings, financed cash flows, terminal risk and policy updates\\
\bottomrule
\end{tabular}
\end{center}

For an overnight risk report, the sensitivity runner can read a canonical
portfolio request and write a summary together with one result per row:
\par\noindent
\begin{minipage}{\linewidth}
\begin{lstlisting}[basicstyle=\ttfamily\normalsize,columns=fullflexible,breaklines=true,frame=single,framesep=3pt]
aad_bench --volt-input book.json \
  --device cpu --precision f64 --gamma \
  --output risk.json \
  --greeks-output rows.jsonl
\end{lstlisting}
\end{minipage}
\par
The same input can be submitted to a GPU-enabled build with the chosen device
and precision. A planning request checks the calculation and estimates its
working memory before opening the device. Larger books can be processed in
chunks while retaining ordered contract results and stable portfolio sums.

The live option service accepts a complete quote-and-position frame through
\texttt{POST /api/quotes}. It calculates bid, midpoint and ask implied
volatilities, fits a surface, updates the book's Greeks and revalues declared
spot and volatility scenarios. A caller can read the latest result through
\texttt{GET /api/snapshot} or subscribe to updates through
\texttt{GET /api/stream}. Hypercube publishes the associated risk snapshot,
which the web and terminal viewers consume. The current live analytics route
runs on CPU; the tensor pricing and AAD interfaces provide the separate
CPU/GPU calculation route shown in Figure~\ref{fig:execution}.

A hedging application can also train a policy from complete scenario batches.
Jormungandr provides the control API, while an optional persistent Volt/Burn
worker retains the policy, optimizer and scenario tensors~\cite{jormungandr}.
The policy chooses stock and option holdings. The worker carries those
holdings through transaction costs, cash funding, dividend income and terminal
liquidation, then differentiates the final risk objective with respect to the
policy weights. Each update has an expected model version, and registered
dataset partitions keep training, validation and test ranges separate.
This learner interface supports the existing continuous MLP policy with
previous-holdings feedback. The companion hedging article develops the
learning problem and its experimental evidence.

\section{Contract values and portfolio exposure}
A European call gives its holder the right to buy at the agreed strike on
expiry; a put gives the right to sell. Their expiry payoffs are $(S_T-K)^+$
and $(K-S_T)^+$, where $x^+=\max(x,0)$. A pricing model estimates the current
value of one option per underlying unit. The
book tells us how many units are owned. For contract $i$, write its model
price as $C_i$, signed quantity as $n_i$, and contract multiplier as $m_i$.
The position weight is $w_i=n_i m_i$, giving portfolio value
\begin{equation}
 V=\sum_{i=1}^{N} w_i C_i.\label{eq:portfolio}
\end{equation}
That same weight applies to each of the contract's per-unit Greeks. A short
position contributes a negative weight. Its contract still describes the same
payoff; the sign records which side of that payoff the book owns.

Take a European call and put on the same stock with
\[
 S=K=100,\quad T=0.5,\quad \sigma=0.20,\quad r=0.03,\quad q=0.01.
\]
Here $S$ is spot, $K$ the strike, $T$ the remaining years, $\sigma$ annual
volatility, $r$ the continuously compounded rate and $q$ the dividend yield.
The example portfolio holds ten calls and is short four puts, each with
multiplier 100. The Black--Scholes formula gives
\begin{center}
\begin{tabular}{@{}lrr@{}}
\toprule
Per underlying unit & Call & Put\\
\midrule
Price &6.090127&5.100073\\
Delta &0.553457&$-0.441555$\\
Vega per unit $\sigma$ &27.789495&27.789495\\
Gamma &0.0277895&0.0277895\\
\bottomrule
\end{tabular}
\end{center}
The weights are 1000 and $-400$, so
\begin{align*}
 V&=1000C-400P=4050.097918,\\
 \Delta_V&=1000\Delta_C-400\Delta_P=730.079337,\\
 \nu_V&=600\times27.789495=16673.696849,\\
 \Gamma_V&=600\times0.0277895=16.673697.
\end{align*}
\begin{center}
\begin{minipage}{\linewidth}
\centering
\begin{tikzpicture}[>=Latex,font=\normalsize,
 box/.style={draw=darkblue,rounded corners=2pt,align=center,inner sep=4pt}]
 \node[box,text width=2.8cm] (c) at(-1.65,0) {10 calls $\times$ 100\\$w_C=1000$};
 \node[box,text width=2.8cm] (p) at(1.65,0) {$-4$ puts $\times$ 100\\$w_P=-400$};
 \node[box,text width=5.5cm] (v) at(0,-1.35)
 {Value $=1000C-400P$\\Delta $=1000\Delta_C-400\Delta_P$};
 \draw[->] (c.south)--(v.north west);
 \draw[->] (p.south)--(v.north east);
\end{tikzpicture}
\captionof{figure}{Quantities and multipliers give signed underlying units.
The same weights carry each contract's price and Greeks into the portfolio.}
\end{minipage}
\end{center}
A volatility move from 20\% to 21\% is a change of 0.01 in
$\sigma$, making the first-order vega effect about 166.737 currency units.
For a one-unit spot rise, delta and gamma give
\[
 \Delta V\approx\Delta_V\Delta S+\tfrac12\Gamma_V(\Delta S)^2
 =738.416185.
\]
In this calculation we hold volatility and the other model inputs fixed.
A surface that moves with spot defines a further sensitivity, developed in
the companion article on live option risk.

Rho is $\partial V/\partial r$,
so a one-basis-point rate move multiplies it by $10^{-4}$. Calendar theta is
$-\partial V/\partial T$, because remaining time falls as the calendar advances.
The adapter uses Actual/365 Fixed for maturity; its one-day linear theta
approximation divides the annual value by 365. A book spanning several
currencies would also need an explicit FX conversion rule before aggregation.

\section{Automatic differentiation for Greeks}
The Black--Scholes price of a European call is
\begin{align}
 C&=Se^{-qT}\Phi(d_1)-Ke^{-rT}\Phi(d_2),\label{eq:bs}\\
 d_1&=\frac{\log(S/K)+(r-q+\sigma^2/2)T}{\sigma\sqrt T},\nonumber\\
 d_2&=d_1-\sigma\sqrt T.\nonumber
\end{align}
The normal cumulative distribution is $\Phi$, with density $\phi$.
The put follows from $P=C-Se^{-qT}+Ke^{-rT}$~\cite{blackscholes}.
For positive spot, strike, maturity and volatility, we can differentiate
through the logarithm, square root, exponentials and normal probabilities.

Automatic differentiation applies the chain rule to these operations as the
program runs~\cite{baydin}. Reverse mode starts from the final scalar and
works back through the intermediate values. The sensitivity of the result to
an intermediate value is called its adjoint. If $z=f(x,y)$ and its adjoint is
$\bar z$, the local reverse step adds
\[
 \bar x\mathrel{+}=\bar z f_x,\qquad
 \bar y\mathrel{+}=\bar z f_y.
\]
The additions allow an input used in several places to receive every
contribution. For example, with $a=\log S$, $b=a^2$ and $V=wb$, the reverse
calculation starts at $\bar V=1$ and gives
\[
 \bar b=w,\qquad \bar a=2a\bar b,\qquad \bar S=\bar a/S.
\]
A larger pricing program repeats these same small steps. It retains or
recomputes the intermediate values needed by each derivative rule, which
explains the additional memory used during a reverse calculation.

For Black--Scholes the independent checks include
\[
 \Delta_C=e^{-qT}\Phi(d_1),\quad
 \nu_C=Se^{-qT}\phi(d_1)\sqrt T,\quad
 \Gamma_C=\frac{e^{-qT}\phi(d_1)}{S\sigma\sqrt T}.
\]
We compare each computed Black--Scholes Greek with these analytic formulas,
which establish its units and provide an independent accuracy check. When a
pricing program contains more steps, we can apply the same local derivative
rules throughout that program. Giles and Glasserman develop this approach
for Monte Carlo risk~\cite{gilesglasserman}.

For stocks exposed to common market factors, let $x_j$ be a log-price shock
and $B_{ij}$ its loading on asset $i$. Define
\begin{equation}
 S_i(x)=S_{i0}\exp\!\left(\sum_{j=1}^{D}B_{ij}x_j\right),\qquad
 V(x)=\sum_i w_iC_i(S_i(x)).\label{eq:factors}
\end{equation}
Reversing that calculation gives all the factor sensitivities:
\begin{equation}
 \frac{\partial V}{\partial x_j}
 =\sum_i B_{ij}w_i\Delta_i S_i.\label{eq:factorgradient}
\end{equation}
A unit loading with $x_j=0.01$ produces approximately a 1\% spot move. The
loadings describe this risk experiment; a statistical scenario model would
also supply a distribution for the factors and their dependence.

\begin{center}
\begin{minipage}{\linewidth}
\centering
\begin{tikzpicture}[>=Latex,font=\normalsize,node distance=5mm,
 box/.style={draw=darkblue,rounded corners=2pt,align=center,text width=4.1cm,inner sep=4pt}]
 \node[box] (x) {Factor shocks $x$};
 \node[box,below=of x] (s) {Asset spots $S_0\exp(Bx)$};
 \node[box,below=of s] (c) {Option prices $C_i$};
 \node[box,below=of c] (v) {Portfolio value $\sum_i w_iC_i$};
 \draw[->] (x)--(s);\draw[->] (s)--(c);\draw[->] (c)--(v);
 \draw[->,darkgreen] (v.east)--++(6mm,0)|-(x.east);
 \node[rotate=90,darkgreen] at(3.05,-1.5) {Reverse sensitivities};
\end{tikzpicture}
\captionof{figure}{The forward calculation builds one portfolio value. The
reverse calculation follows its dependence on every factor.}
\end{minipage}
\end{center}
Central differences would reprice the book at $x+he_j$ and $x-he_j$ for each
factor, requiring $2D$ prices. Reverse mode shares the pricing work across all
$D$ derivatives. Forming the dense projection $Bx$ and its reverse still costs
on the order of $ND$, and the program returns $D$ results. This is the useful
computational saving: expensive pricing work is reused across its inputs.

For Black--Scholes spot gamma, Volt carries three values through each operation: its value, first derivative and
second derivative in the spot direction. Seeding spot with $(S,1,0)$ and
constants with $(c,0,0)$ gives, for example,
\[
 (fg)''=f''g+2f'g'+fg'',\qquad
 (\log f)''=f''/f-(f'/f)^2.
\]
Applying these rules through Equation~\eqref{eq:bs} gives spot gamma directly.
A change of coordinates adds its own curvature. In a one-factor log-price
experiment, the corresponding second derivative is
\[
 \frac{\partial^2V}{\partial x^2}
 =\sum_i w_iB_i^2(\Gamma_i S_i^2+\Delta_i S_i).
\]
The delta term follows from differentiating the exponential spot map twice.


A simulated payoff raises a different issue. Under risk-neutral geometric Brownian motion
(GBM), a call price is the discounted expectation of $(S_T-K)^+$, with
\[
 S_T=S\exp\!\left((r-q-\tfrac12\sigma^2)T+\sigma\sqrt T Z\right),
 \qquad Z\sim N(0,1).
\]
For each sampled path, the first derivative propagates through $S_T$ and the
payoff. The payoff's slope changes at the strike, however, and its indicator
is constant on either side. A second pathwise differentiation misses the
probability mass crossing that strike as spot changes.

We account for that movement with a likelihood-ratio estimator, which
weights each discounted payoff by
\begin{equation}
 W_\Gamma(Z)=\frac{Z^2-1}{S^2\sigma^2T}
             -\frac{Z}{S^2\sigma\sqrt T}.\label{eq:lrgamma}
\end{equation}
This estimates the second derivative of the expected price~\cite{broadie}.
Its sampling error is measured alongside the estimate. When using paired
normal draws $(Z,-Z)$, each pair supplies one independent observation.
For a portfolio, common draws also preserve covariance between its estimated
position risks. A multi-asset market model supplies the economic dependence
between assets.

\section{GPU execution and memory}
The implementation uses Burn for tensors and automatic differentiation, and
CubeCL for accelerator execution~\cite{burn,cubecl}. A tensor is an array with
a shape, a number format and an execution device. These choices are explicit
in a Volt session. The CPU path supports 32-bit and 64-bit floating point
(f32 and f64), Metal supports the tested f32 path, and CUDA supports both.
The session checks the requested combination before executing it.

Hypercube handles the publication side of the live application. It groups
calculated prices and risks with their observation times and portfolio version
into a coherent snapshot. The current live surface application calculates its
IVs, fitted surface and Greeks through the CPU analytics route. Batched pricing
and AAD requests use Burn's tensor route, which can select CPU or GPU execution.
Figure~\ref{fig:execution} shows where those routes sit in the application.



An application array is packed into the selected number format and uploaded
into a device tensor. Kernel operations then use those tensors, and a readback
returns the requested results to the host. Matrix shape is part of the request:
an $m\times k$ array multiplied by a $k\times n$ array produces an
$m\times n$ result. Rectangular and odd-sized fixtures exercise the layouts
and partially filled work tiles that arise in larger financial calculations.

On Apple Silicon, CPU and GPU share physical memory. Metal still assigns a
storage mode to each resource and coordinates access to it~\cite{apple}.
The current tensor API handles that allocation and upload. The application
retains its host inputs while the runtime manages the device resources. A discrete NVIDIA
GPU adds a transfer between host memory and device memory. Reusing inputs on
the device spreads this exchange cost over more calculations~\cite{nvidia}.
That reuse is called residency.

For one request, the elapsed time has several parts:
\begin{align*}
 t_{\rm job}={}&t_{\rm setup}+t_{\rm pack}+t_{\rm upload}
 +t_{\rm calculation}\\
 &+t_{\rm synchronization}+t_{\rm readback}+t_{\rm aggregation}.
\end{align*}
The first request may also initialize a device context and compile kernels.
Enqueuing GPU work can return while the device is still calculating, so a
completed-work timing includes the required synchronization.

Volt's persistent valuation route uploads its inputs once and can apply
supported uniform spot changes on the device. Each current AAD request builds
its own derivative graph. These routes therefore have different reuse and
storage costs, even when they start from the same option rows. The measurements
below show the consequence: resident valuation benefits from CUDA at a tested
size where a complete price-and-Greeks request remains faster on CPU.

For the current f32 first-order AAD workload, the memory planner allows 256 scalar graph slots, 512 bytes
of staging and 48 bytes of returned values per row. At 65,536 rows these terms
are 64 MiB, 32 MiB and 3 MiB, totaling 99 MiB before the other declared reserves
and headroom. The graph term becomes 128 MiB at f64; a gamma calculation receives
its own additional allowance. These estimates let the planner choose a batch
size before allocating the work.

If that work exceeds the allowance, the calculation can proceed in row chunks.
A portfolio sum combines the chunks in source order, and the shared-factor
example combines each chunk's $D$-component gradient in the same way. Smaller
chunks reduce live memory and repeat more setup, calculation and readback work.
We measured this tradeoff directly: for 1,024 positions and 256 factors on
an Apple M3 using Metal f32, increasing the chunk from 32 to 1,024 rows raised sampled live
allocation from about 38 KB to 1.10 MB while reducing elapsed time from roughly
112 ms to 6.1 ms. The aggregate gradient passed the same accuracy checks.
These observations came from a host running other work.

The 99 MiB allowance above estimates array storage and derivative workspace.
Allocator samples measure storage at a particular point during execution;
process memory also includes host and runtime state. Caller-owned market data, compiler storage and other
applications enter the deployment's memory plan separately.

\section{Heston pricing and numerical validation}
Heston models variance as a process that evolves alongside the stock price~\cite{heston}:
\begin{align*}
 dS_t&=(r-q)S_t\,dt+\sqrt{v_t}S_t\,dW_t^S,\\
 dv_t&=\kappa(\vartheta-v_t)\,dt+\xi\sqrt{v_t}\,dW_t^v,\\
 d\langle W^S,W^v\rangle_t&=\rho\,dt.
\end{align*}
Variance $v_t$ returns towards its long-run level $\vartheta$ at speed
$\kappa$. The parameter $\xi$ controls its random fluctuations, while $\rho$
links those fluctuations to stock moves through the two Brownian drivers
$W^S$ and $W^v$. The model's characteristic function,
the expectation of a complex exponential of log price, lets us calculate
European option values through Fourier integrals.

The GPU evaluates those integrals on a fixed numerical grid. A
Black--Scholes control variate improves the arithmetic: we add its known
price and integrate the difference between the two characteristic functions.
Heston's dynamics continue to determine the result. Differentiating this
calculation supplies sensitivities to the pricing inputs and variance
parameters; Fourier spot derivatives also supply gamma.

Heston's initial variance is
$v_0$, so $C_{v_0}$ measures price per variance unit. If we choose to express
$v_0=\sigma_0^2$, the chain rule gives
\[
 \frac{\partial C}{\partial\sigma_0}=2\sigma_0 C_{v_0}.
\]
The companion
hedging paper uses these model sensitivities to construct its stock and
variance-instrument comparison strategy.

Analytic Black--Scholes Greeks test the smooth pricing graph. Central
repricing and the analytic factor expression in Equation~\eqref{eq:factorgradient}
verify shared-factor sensitivities. Monte Carlo estimates are compared with
analytic prices and Greeks using their sampling uncertainty. Heston has a
separate CPU-f64 reference that adds integration points where the estimated error is larger,
with explicit evaluation limits and a tail estimate within its stated parameter domain~\cite{quadpack}.

In central differences, a smaller bump reduces truncation error but increases subtraction roundoff,
whose contribution scales roughly as $\epsilon |V|/h$. We compare several
bumps. For a Fourier integral, extra nodes improve accuracy within the chosen
interval; extending the cutoff includes more of the tail and introduces more
oscillations to resolve. We vary node count and cutoff separately. This is
particularly useful near zero variance and short maturity, where the tail can
be demanding.

An f64 derivative rule used a lower-precision reciprocal within the normal-CDF
calculation. Direct Greek checks exposed the error, and expressing the constant
operation through its reciprocal restored the expected accuracy. On CUDA,
native special functions through the C++ backend and NVIDIA's runtime CUDA
compiler (NVRTC) supplied the precision required by these checks.

Dense matrix fixtures exposed reduced-precision multiplication in the selected
CUDA route, despite the arrays being stored as f32. The supported CUDA matrix route uses
strict multiplication, and small differentiable policy layers use bounded
products and reductions. For portfolio totals, ordered compensated f64 sums
preserve small residual exposures when large positive and negative positions
nearly offset.

The numerical checks cover prices, first derivatives, spot gamma and shared
factor sensitivities on the supported CPU, Metal and CUDA routes. Each route
is checked against the corresponding independent formula, numerical reference
or statistical estimate. The shared-factor checks cover 256 positions with
1, 4, 16, 64 and 256 inputs, using both analytic derivatives and central
repricing. The supported precision is explicit: Metal uses f32,
while CPU and CUDA provide f32 and f64.

\section{Choosing an execution route}
We measured the supplied calculation tools on a system with an NVIDIA RTX 4090
and an AMD Ryzen Threadripper PRO 7985WX processor running Linux~\cite{volt}. The CPU results
use one worker. The machine was under low load during this measurement series,
and warm timings follow device initialization. GPU timings include the
synchronization needed to complete the requested calculation.

For matrix multiplication, input
hashing and independent validation are measured separately. The complete Greeks
request includes setup, transfer, differentiation, gamma, readback and
aggregation; its analytic validation follows the timed calculation. The largest
matrix uses structured inputs with a closed-form reference answer,
while the device performs the full dense multiplication. The smaller matrices
use independent dense inputs.

\begin{center}
\begin{minipage}{\linewidth}
\centering
\begin{tabular}{@{}p{4.7cm}rr@{}}
\toprule
Warm workload & \multicolumn{2}{c}{Time in milliseconds}\\
 & CPU & CUDA\\
\midrule
Matrix $64^3$, f32 &0.00493&0.04196\\
Matrix $256^3$, f32 &0.23943&0.06656\\
Matrix $1024^3$, f64 &35.02391&7.11147\\
Resident valuation, 4096 rows, f32 &0.11275&0.41250\\
Resident valuation, 65536 rows, f32 &1.74625&0.67028\\
Price, first Greeks and gamma, 65536 rows, f32 &24.24139&27.66377\\
\bottomrule
\end{tabular}
\captionof{table}{Selected measured medians under the stated execution contract.
One millisecond (ms) is 0.001 second. The last row uses a requested chunk cap
of 65,536 on both backends. The matrix,
resident-valuation and complete-Greeks rows each retain their timing scope.}
\end{minipage}
\end{center}

CPU handles the small matrix and the small resident valuation faster. CUDA is faster at the larger displayed sizes,
including about a 2.6-fold advantage for the 65,536-row resident valuation.
The full Greeks request takes 24.24 ms on CPU and 27.66 ms on CUDA. The table
therefore supports choosing execution by request type and data reuse. An
application can keep CPU as its default for a modest complete Greek report
and choose the GPU for a measured larger resident workload.

\end{SNcolumns}
\clearpage
\begin{SNcolumns}
\begin{thebibliography}{99}
\interlinepenalty=10000
\bibitem{blackscholes} F. Black and M. Scholes.
The pricing of options and corporate liabilities.
\emph{Journal of Political Economy}, 81(3):637--654, 1973.
\url{https://doi.org/10.1086/260062}.
\bibitem{baydin} A. G. Baydin, B. A. Pearlmutter, A. A. Radul and J. M. Siskind.
Automatic differentiation in machine learning: a survey.
\emph{Journal of Machine Learning Research}, 18(153):1--43, 2018.
\url{https://jmlr.org/papers/v18/17-468.html}.
\bibitem{gilesglasserman} M. Giles and P. Glasserman.
Smoking adjoints: fast Monte Carlo Greeks.
\emph{Risk}, 19(1):88--92, 2006.
Author's implementation and paper resources:
\url{https://people.maths.ox.ac.uk/~gilesm/codes/libor_AD/}.
\bibitem{broadie} M. Broadie and P. Glasserman.
Estimating security price derivatives using simulation.
\emph{Management Science}, 42(2):269--285, 1996.
\url{https://doi.org/10.1287/mnsc.42.2.269}.
\bibitem{heston} S. L. Heston.
A closed-form solution for options with stochastic volatility with applications
to bond and currency options.
\emph{Review of Financial Studies}, 6(2):327--343, 1993.
\url{https://doi.org/10.1093/rfs/6.2.327}.
\bibitem{quadpack} R. Piessens, E. de Doncker-Kapenga, C. W. \"Uberhuber and
D. K. Kahaner. \emph{QUADPACK: A Subroutine Package for Automatic Integration}.
Springer, 1983. Original routine \texttt{DQK15}:
\url{https://netlib.org/quadpack/dqk15.f}.
\bibitem{burn} Tracel. \emph{Burn}, version 0.22.1. Source and documentation.
\url{https://github.com/tracel-ai/burn/tree/v0.22.1}.
\bibitem{cubecl} Tracel. \emph{CubeCL}, version 0.11.3. Source and documentation.
\url{https://github.com/tracel-ai/cubecl/tree/v0.11.3}.
\bibitem{apple} Apple. Choosing a resource storage mode for Apple GPUs.
\emph{Metal developer documentation}.
\url{https://developer.apple.com/documentation/metal/choosing-a-resource-storage-mode-for-apple-gpus}.
\bibitem{nvidia} NVIDIA. Data transfer between host and device.
\emph{CUDA C++ Best Practices Guide}, section on data transfer.
\url{https://docs.nvidia.com/cuda/cuda-c-best-practices-guide/index.html#data-transfer-between-host-and-device}.
\bibitem{volt} StrategyNet. \emph{Volt: option pricing, risk and hedging APIs}.
Source and documentation.
\url{https://github.com/ingvor-restrat/volt}.
\bibitem{jormungandr} StrategyNet. \emph{Jormungandr: learner services and native application support}. Source and documentation.
\url{https://github.com/ingvor-restrat/jormungandr}.
\end{thebibliography}
\end{SNcolumns}
\end{document}
